\INSRAbstract{
Integrative Neuro-Somatic Recalibration (INSR) is proposed as a state-adaptive decision architecture for sequencing, pausing, resuming, and combining interventions in complex trauma care. Rather than defining another fixed treatment protocol, INSR formalizes clinical decisions about regulation, processing readiness, dissociation, mentalization, self-state conflict, and safety. The framework integrates psychotherapeutic principles with optional physiological and neurotechnological measurement layers while retaining clinician oversight. This position paper outlines the conceptual rationale, proposed gating architecture, translational hypotheses, safety constraints, and an empirical roadmap for feasibility testing. INSR remains a research framework under development and is not presented as a validated clinical intervention.
}

\INSRKeywords
